% Proposed model interface. No experimental response or anatomical transfer.
\begin{tikzpicture}[x=1mm,y=1mm,font=\fontsize{12}{14}\selectfont,
  line width=.55pt,>=Latex]
\path[use as bounding box] (0,0) rectangle (160,83);
\node[anchor=west,font=\bfseries] at (1,80) {Предобучение на Drosophila};
\node[anchor=west,font=\bfseries] at (1,38) {Целевая модель человеческого сигнала};
% Source connectivity, observed trajectories and learned initial parameters.
\coordinate (a) at (5,62); \coordinate (b) at (15,70);
\coordinate (c) at (17,54); \coordinate (d) at (28,64);
\coordinate (e) at (38,55);
\foreach \a/\b in {a/b,a/c,b/d,c/d,d/e,c/e} \draw[->,gray] (\a)--(\b);
\foreach \p in {a,b,c,d,e} \fill (\p) circle (.9);
\node at (22,47) {$G_F,\theta_F$};
\draw[->] (44,63)--(53,63);
\node at (48,72) {$F,Q_F$};
\draw[->] (56,52)--(98,52) node[below left] {$t$};
\draw[->] (56,52)--(56,71) node[above right] {$o_F$};
\draw[thick] plot[smooth] coordinates {(58,57)(62,58)(66,68)(70,57)(75,61)(80,66)(85,55)(91,62)(97,59)};
\node[align=center] at (79,48) {$P_F,\ E_\omega,\ D_F$};
\draw[->] (102,63)--(113,63);
\foreach \x/\h in {118/7,123/13,128/5,133/10,138/16,143/8} {
  \draw[gray] (\x,53)--(\x,53+\h);
  \fill (\x,53+\h) circle (.7);
}
\node[anchor=west] at (145,62) {$\omega_0$};
\draw[dashed,->] (132,49)--(132,34)--(29,34)--(29,31);
% Target temporal representation is distinct from physical axonal currents.
\draw[->] (3,10)--(44,10) node[below left] {$t$};
\draw[thick] plot coordinates {(4,13)(8,13)(8,24)(11,24)(11,13)(19,13)(19,24)(22,24)(22,13)(32,13)(32,24)(35,24)(35,13)(43,13)};
\node at (25,29) {$u_H,c_H:\ P_H,E_\omega$};
\draw[->] (47,20)--(57,20) node[midway,below] {$B_\alpha$};
\node at (52,30) {$u_H$};
\draw[->] (52,27)--(52,21);
\foreach \y/\k in {12/0,18/1,24/2} {
  \draw[gray] (61,\y)--(97,\y);
  \foreach \x in {64,74,84,94} \fill (\x,\y) circle (.6);
  \draw[thick] (65+\k*5,\y)--(67+\k*5,\y+5)--(69+\k*5,\y-2)--(71+\k*5,\y);
}
\node[align=center] at (79,3) {$r_H$: токи сегментов, А};
\draw[->] (100,20)--(113,20) node[midway,above] {$H_\varphi$};
% Passive recording produces voltage; amplitude depends on stimulus current.
\draw[->] (116,15)--(134,15) node[below left] {$t$};
\draw[->] (116,7)--(116,29) node[above right] {$\hat y_H$};
\draw[thick] plot[smooth] coordinates {(117,15)(120,16)(122,23)(124,8)(127,19)(130,16)(133,15)};
\draw[->] (134,12)--(139,12);
\draw[->] (141,9)--(158,9) node[below left] {$I$};
\draw[->] (141,9)--(141,29) node[above right] {$A$};
\draw[thick] plot[smooth] coordinates {(142,10)(146,10)(149,13)(152,18)(157,24)};
\end{tikzpicture}
